\chapter{Default Theories, Arguments and Autoepistemic Expansions}
\label{sec:p18-default-autoepistemic-argumentation}

Dung's Section~4.1 gives an argument construction for
Reiter's closed default theories. We supply its full
closure, consistency and stable-extension proofs, then
state exactly what the direct autoepistemic translation
preserves \cite[pp.~339--341]{Dung1995}.

\begin{hbdefinition}{def:p18-closed-default-theory}{Closed theories, justifications and classical consequence}
Fix a first-order language with its set \(\mathcal L_0\)
of closed sentences. Write \(\operatorname{Th}(X)\)
for classical deductive closure. We use its finite-proof
property: each consequence has a finite subset of
premises. A set of sentences is consistent when it does not
entail \(\bot\). Sentences are written with double
negations removed. Their normalized contrary
\(\bar\varphi\) is classically equivalent to
\(\neg\varphi\) and satisfies
\(\overline{\bar\varphi}=\varphi\).

A closed default is
\[
                 d=\frac{\alpha:\beta_1,\ldots,\beta_k}{\gamma},
\]
where all displayed formulas are closed and the
justification list is finite. A default theory
\(\mathcal T=(D,W)\) has a set of such defaults
and a set of background sentences, either of
which may be infinite. Let \(J=\operatorname{Jus}(D)\).
Empty justification lists are permitted here;
their effect on consistency is stated explicitly
below. A missing prerequisite means \(\top\).
Justifications are tested individually, not by
requiring their conjunction to be consistent.
They license default steps and are not added
as ordinary premises.
Background consistency and existence of a
consistent Reiter extension are kept as
separate conditions.
\end{hbdefinition}

\begin{hbdefinition}{def:p18-default-closure}{Licensed consequence and Reiter's operator}
For \(B\subseteq J\), set \(C_0(B)=W\) and
\[
 C_{n+1}(B)=\operatorname{Th}(C_n(B))\cup
   \{\gamma_d:d\in D,\ \alpha_d\in C_n(B),\
                              \{\beta_{d,1},\ldots,\beta_{d,k}\}\subseteq B\},
 \qquad C(B)=\bigcup_{n<\omega}C_n(B).
\]
The unconditional background closure is \(C(\varnothing)\).
For a deductively closed candidate theory \(E\), put
\[
 B(E)=\{j\in J:E\cup\{j\}\text{ is consistent}\}
     =\{j\in J:\bar j\notin E\},\qquad
                     \Gamma_{\mathcal T}(E)=C(B(E)).
\]
A Reiter extension is \(E=\Gamma_{\mathcal T}(E)\).
Restricting candidates to closed theories loses
no extension: every licensed-closure value is closed.
This is the original final-candidate consistency test:
the test uses \(E\), while prerequisites must be
generated in the increasing construction from \(W\)
\cite[definition~1, theorem~2.1, pp.~88--90]{Reiter1980}.

A finite support \(K\subseteq J\) supports \(\varphi\)
when there is a finite derivation ending in \(\varphi\),
using members of \(W\), classical consequences of
earlier members, and defaults whose prerequisites
occur earlier and whose justifications belong to \(K\).
Unused elements of \(K\) are allowed.
\end{hbdefinition}

\begin{hblemma}{thm:p18-default-finite-support}{proved=annex:p18.default.finite.support}{The complete finite-support and closure interface}
\(C(B)\) is the least classically closed set containing
\(W\) and closed under the defaults licensed by \(B\).
It is monotone in \(B\), and
\[
 \varphi\in C(B)\quad\Longleftrightarrow\quad
 \text{some finite }K\subseteq B\text{ supports }\varphi.
\]
These conclusions allow infinite \(D,W,B\).
\end{hblemma}

\begin{hbtheorem}{thm:p18-default-extension-properties}{proved=annex:p18.default.extension.properties}{Generating defaults, consistency and incomparable extensions}
Every Reiter extension \(E\) satisfies
\[
 E=\operatorname{Th}\!\left(W\cup
       \{\gamma_d:d\in D,\ \alpha_d\in E,\
                              \bar\beta_{d,i}\notin E\text{ for all }i\}\right).
\]
The map \(\Gamma_{\mathcal T}\) reverses inclusion
on closed theories. Consequently two comparable
Reiter extensions are equal.

An inconsistent extension exists exactly when
\(C(\varnothing)\) is inconsistent; in that
case the sole extension is \(\mathcal L_0\).
If every default has at least one justification,
\(C(\varnothing)=\operatorname{Th}(W)\), recovering
the original background-consistency statement.
Consistent background alone does not guarantee
any extension.
\end{hbtheorem}

\section{The literal default-theory argument construction}

\begin{hbdefinition}{def:p18-default-argument-graph}{Arguments, attacks and the two maps}
The argument set of \(\mathcal A(\mathcal T)\) is
\[
          \{(K,\varphi):K\subseteq J\text{ is finite and supports }\varphi\}.
\]
The attack clause is exactly contrary membership:
\[
                (K,\varphi)\to(K',\psi)
                                \quad\Longleftrightarrow\quad\bar\varphi\in K'.
\]
The normalized contrary convention makes this
equivalently an argument for \(\bar j\) attacking
a support containing \(j\).
For an argument set \(S\), define
\[
 \operatorname{flat}(S)=\{\varphi:\exists K\ (K,\varphi)\in S\},
 \qquad H(S)=\bigcup_{(K,\varphi)\in S}K.
\]
For a closed theory \(E\), define
\[
 \operatorname{arg}(E)=
       \{(K,\varphi)\in\mathcal A(\mathcal T):K\subseteq B(E)\}.
\]
Membership in a framework means membership in
its argument set. Reuse the characteristic,
admissible, complete, preferred and stable
definitions of the preceding chapters.
\end{hbdefinition}

\begin{hblemma}{thm:p18-default-admissible-consistency}{proved=annex:p18.default.admissible.consistency}{The full admissible-support consistency lemma}
For every admissible \(S\),
\[
          C(H(S))\text{ is consistent}
                     \quad\Longleftrightarrow\quad C(\varnothing)\text{ is consistent}.
\]
In particular \(\operatorname{flat}(S)\) is
consistent when that unconditional closure is.
If \(S\) is complete, its flattening is
classically closed and contains \(W\).
With nonempty justification lists the first
equivalence is Dung's Lemma~40 in its
background-consistency form.
\end{hblemma}

\begin{hblemma}{thm:p18-default-flattening}{proved=annex:p18.default.flattening}{The complete flattening identity and stable closure}
For every closed theory \(E\),
\[
               \operatorname{flat}(\operatorname{arg}(E))
                                     =\Gamma_{\mathcal T}(E).
\]
For every stable argument extension \(S\),
its flattening is closed and
\[
                       S=\operatorname{arg}(\operatorname{flat}(S)).
\]
The statements include inconsistent theories.
\end{hblemma}

\begin{hbtheorem}{thm:p18-default-stable-bijection}{proved=annex:p18.default.stable.bijection}{Both directions of Dung's default correspondence}
The maps \(\operatorname{arg}\) and
\(\operatorname{flat}\) restrict to inverse
bijections between Reiter extensions of
\(\mathcal T\) and stable extensions of
\(\mathcal A(\mathcal T)\).
Each corresponding stable extension is
preferred. Preferred extensions always
exist, even when Reiter and stable
extensions do not.
\end{hbtheorem}

\begin{hbdefinition}{def:p18-default-examples}{Explicit default theories}
In the following propositional cases all
named atoms are independent and \(p,q,r\)
are distinct:
\[
\begin{aligned}
 \mathcal T_{\rm paradox}
   &=\left(\left\{\frac{\top:\neg p}{p}\right\},\{q\}\right),\\
 \mathcal T_{\rm separate}
   &=\left(\left\{\frac{\top:p,\neg p}{r}\right\},\varnothing\right),\\
 \mathcal T_{\rm self}
   &=\left(\left\{\frac{p:p}{p}\right\},\varnothing\right),\\
 \mathcal T_{\rm inconsistent}
   &=(\varnothing,\{p,\neg p\}),\\
 \mathcal T_{\rm unconditional}
   &=\left(\left\{\frac{\top:}{\bot}\right\},\varnothing\right).
\end{aligned}
\]
\end{hbdefinition}

\begin{hbtheorem}{thm:p18-default-examples}{proved=annex:p18.default.examples}{All separating default and argument calculations}
\(\mathcal T_{\rm paradox}\) has no Reiter
extension and no stable argument extension.
Its unique preferred and grounded argument
set consists of all empty-support arguments;
its flattening is \(\operatorname{Th}(\{q\})\).
Thus the background conclusion \(q\) remains
in that preferred theory despite absence
of a stable-extension family.

\(\mathcal T_{\rm separate}\) has the unique
extension \(\operatorname{Th}(\{r\})\), although
its two justifications are jointly inconsistent.
Its entire argument graph is unattacked
and is its unique stable extension.
\(\mathcal T_{\rm self}\) has the unique
Reiter extension \(\operatorname{Th}(\varnothing)\).
Each of the last two theories has the
sole inconsistent extension \(\mathcal L_0\).
The unconditional case has consistent
\(W\) but inconsistent \(C(\varnothing)\).
\end{hbtheorem}

\begin{hbmethod}{def:p18-default-finite-method}{Finite propositional extension enumeration}
For finite propositional \(D,W\), enumerate
subsets \(G\subseteq D\). Represent the
candidate \(E\) by the finite basis
\(W\cup\{\gamma_d:d\in G\}\).
Test each justification separately against
that candidate, then construct its
\(\Gamma_{\mathcal T}(E)\) by repeatedly
adding licensed conclusions whose
prerequisites follow classically.
Return the candidate exactly when
the two finite bases generate the
same theory, suppressing repetitions.
\end{hbmethod}

\begin{hbtheorem}{thm:p18-default-finite-method}{proved=annex:p18.default.finite.method}{The exact finite algorithm domain}
The procedure terminates and returns exactly
all Reiter extensions. Classical tests can
be performed over the finitely many truth
assignments of the input atoms. The
licensed-conclusion loop has at most
\(|D|\) strict enlargements of its
conclusion list. This is a finite
propositional procedure; the infinite
closed-theory correspondence does not
assume those tests decidable.
\end{hbtheorem}

\section{Ordinary autoepistemic expansions}

\begin{hbdefinition}{def:p18-ael-language}{Closed modalities and the original expansion equation}
Extend \(\mathcal L_0\) by \(L\theta\) only
when \(\theta\) is a closed sentence of
the extended language \(\mathcal L_L\).
Nested modalities are permitted; quantifying
into a modality is not. In ordinary
classical consequence, each closed modal
atom \(L\theta\) is treated as an
independent nullary predicate.
Write \(\operatorname{Cn}_L\) for this
consequence operation.
We use classical first-order soundness
and completeness to identify deductive
consequence with truth in all these
ordinary structures.

For \(T\subseteq\mathcal L_L\), put
\[
 {\cal B}(T)=\{L\theta:\theta\in T\}
                          \cup\{\neg L\theta:\theta\notin T\}.
\]
An ordinary stable expansion of a base
\(A\subseteq\mathcal L_L\) is
\[
                       T=\operatorname{Cn}_L(A\cup{\cal B}(T)).
\]
This is Moore's closure/introspection
equation; the closed-modal first-order
language follows Konolige's explicitly
non-quantifying-in extension
\cite[pp.~274--277]{Moore1983Autoepistemic}
\cite[sec.~2.1]{Konolige1987DefaultAutoepistemic}.
\end{hbdefinition}

\begin{hbdefinition}{def:p18-ael-kernel-completion}{The explicit completion of an objective kernel}
Given a classically closed \(E\subseteq\mathcal L_0\),
define membership in \(\operatorname{St}(E)\)
by induction on modal depth.
For an objective sentence, membership means
membership in \(E\). For a general
sentence \(\theta\), replace each outer
modal atom \(L\psi\) by \(\top\) if
\(\psi\in\operatorname{St}(E)\) and
by \(\bot\) otherwise. Its interior
\(\psi\) has smaller modal depth.
Call the resulting objective sentence
\(\theta^E\). Put
\[
                    \theta\in\operatorname{St}(E)
                                \quad\Longleftrightarrow\quad E\models\theta^E.
\]
If \(E=\mathcal L_0\), this defines
\(\operatorname{St}(E)=\mathcal L_L\).
\end{hbdefinition}

\begin{hbtheorem}{thm:p18-ael-kernel}{proved=annex:p18.ael.kernel}{Kernel completion and reduction of every modal constructor}
\(\operatorname{St}(E)\) is classically
closed, contains its introspective
assignment, has objective part \(E\),
and satisfies
\[
 \operatorname{St}(E)=
                   \operatorname{Cn}_L(E\cup{\cal B}(\operatorname{St}(E))).
\]
It is the unique classically closed
introspective theory with objective
part \(E\). In particular every
stable expansion is determined by
its objective kernel. The construction
handles every finite nesting depth,
Boolean operation and permitted
first-order quantifier.
\end{hbtheorem}

\begin{hbdefinition}{def:p18-default-ael-transform}{The direct transform and its weak fixed-point operator}
The direct modal transform is
\[
 \tau(\mathcal T)=W\cup
 \left\{\left(L\alpha_d\land\bigwedge_i\neg L\bar\beta_{d,i}\right)
                                              \to\gamma_d:d\in D\right\}.
\]
For a closed objective theory \(E\), define
\[
 \Delta_{\mathcal T}(E)=
 \operatorname{Th}\!\left(W\cup
 \{\gamma_d:\alpha_d\in E,\ \bar\beta_{d,i}\notin E\text{ for every }i\}\right).
\]
A weak extension here means exactly
\(E=\Delta_{\mathcal T}(E)\).
Its prerequisite test uses the final
candidate directly, unlike Reiter's
least licensed closure.
\end{hbdefinition}

\begin{hbtheorem}{thm:p18-default-ael-weak}{proved=annex:p18.default.ael.weak}{The exact general meaning of the direct translation}
Ordinary stable expansions of \(\tau(\mathcal T)\)
correspond bijectively to the weak extensions
defined above, by taking the objective
kernel and applying \(\operatorname{St}\).
Every Reiter extension is such a weak
extension, and hence gives an ordinary
stable expansion. The converse fails
in general.
\end{hbtheorem}

\begin{hbtheorem}{thm:p18-default-ael-prerequisite-free}{proved=annex:p18.default.ael.prerequisite.free}{The full prerequisite-free bijection}
If every prerequisite is \(\top\), then
\(\Gamma_{\mathcal T}(E)=\Delta_{\mathcal T}(E)\)
for every closed candidate \(E\). Thus
Reiter extensions, ordinary stable expansions
of the direct transform, and stable
argument extensions of Dung's default
graph correspond bijectively, including
the inconsistent case. In the transform
the conjunct \(L\top\) may be omitted.
The result permits infinitely many
closed defaults and background sentences.
\end{hbtheorem}

\begin{hbtheorem}{thm:p18-default-ael-boundary}{proved=annex:p18.default.ael.boundary}{Self-support and the exact remaining translation boundary}
For \(\mathcal T_{\rm self}\), the direct
transform is
\[
                         (Lp\land\neg L\neg p)\to p.
\]
It has exactly two ordinary stable
expansions, with kernels
\(\operatorname{Th}(\varnothing)\) and
\(\operatorname{Th}(\{p\})\).
Only the first kernel is a Reiter
extension. Hence a direct identification
of all ordinary expansions with Reiter
extensions is false.

A translation of arbitrary default
theories in both directions with
ordinary autoepistemic logic, or a
separate strongly grounded expansion
construction, requires its own exact
proof. The prerequisite-free theorem
does not establish that broader target.
\end{hbtheorem}

The full Dung Section~4.1 default argument
correspondence is proved above. The general
autoepistemic translation boundary is
kept explicit. Section~3 games and the complete Lucas construction
are supplied in \cref{sec:p18-game-argumentation,sec:p18-lucas-no-solution}.
Section~5's full meta-interpreter and its specializations are
\cref{sec:p18-argumentation-meta-interpreters}.
The separate Dung 1991 presentation retains its distinct source
development.
